%%%%%%%%%%%%%%%%%%%%%%%%%%%%%%%%%%%%%%%%%%%%%%%%%%%%%%%%%%%%%%%%
%  chapter2.tex  --  Implementacao
%  Jogo ISOLA em Prolog com Algoritmo Alfa-Beta
%  Inteligencia Artificial, ISEL 2025/2026
%%%%%%%%%%%%%%%%%%%%%%%%%%%%%%%%%%%%%%%%%%%%%%%%%%%%%%%%%%%%%%%%

\chapter{Implementa\c{c}\~{a}o}
\label{chap:impl}

Este cap\'{i}tulo descreve a implementa\c{c}\~{a}o do jogo ISOLA em SWI-Prolog.
O projeto est\'{a} organizado em cinco m\'{o}dulos:

\begin{table}[htbp]
  \caption{Estrutura de m\'{o}dulos do projeto}
  \label{tab:modules}
  \centering
  \begin{tabular}{lll}
    \toprule
    \textbf{Ficheiro} & \textbf{Responsabilidade} & \textbf{Predicados principais} \\
    \midrule
    \texttt{board.pl}  & Representa\c{c}\~{a}o e display do tabuleiro
                       & \texttt{initial\_board}, \texttt{get\_cell}, \texttt{set\_cell} \\
    \texttt{rules.pl}  & Regras e valida\c{c}\~{a}o de movimentos
                       & \texttt{valid\_move}, \texttt{apply\_move}, \texttt{game\_over} \\
    \texttt{ai.pl}     & Intelig\^{e}ncia artificial (alfa-beta)
                       & \texttt{best\_move}, \texttt{alphabeta}, \texttt{evaluate} \\
    \texttt{game.pl}   & Interface e ciclo de jogo
                       & \texttt{play}, \texttt{game\_loop}, \texttt{human\_turn} \\
    \texttt{main.pl}   & Ponto de entrada
                       & \texttt{initialization} \\
    \bottomrule
  \end{tabular}
\end{table}

% ===================================================================
\section{Representa\c{c}\~{a}o do Tabuleiro (\texttt{board.pl})}
\label{sec:board}

\subsection{Estrutura de Dados}

O tabuleiro \'{e} representado como uma \textbf{lista de listas} em Prolog.
Cada elemento da lista exterior \'{e} uma linha do tabuleiro; cada elemento
de uma linha \'{e} um \'{a}tomo que indica o conte\'{u}do da casa:

\begin{itemize}
  \item \texttt{'.'} --- casa livre (vazia)
  \item \texttt{x}   --- pe\~{a}o do Jogador~1
  \item \texttt{o}   --- pe\~{a}o do Jogador~2
  \item \texttt{'\#'} --- casa permanentemente removida
\end{itemize}

As posi\c{c}\~{o}es s\~{a}o representadas como o termo composto \texttt{Linha/Coluna}
(indexa\c{c}\~{a}o a partir de zero). Por exemplo, \texttt{2/3} significa linha~2,
coluna~3.

\textbf{Vantagem desta representa\c{c}\~{a}o:} as listas s\~{a}o a estrutura
nativa do Prolog, e predicados como \texttt{nth0/3} (acesso por \'{i}ndice) e
\texttt{maplist/2} (aplica\c{c}\~{a}o funcional) est\~{a}o dispon\'{i}veis na
biblioteca padr\~{a}o, o que simplifica enormemente o c\'{o}digo.

\subsection{Predicados Principais}

O predicado \texttt{initial\_board/2} cria um tabuleiro de tamanho
$N{\times}N$ totalmente preenchido com casas livres:

\begin{lstlisting}[language=Prolog,
    caption={Cria\c{c}\~{a}o do tabuleiro inicial},
    label={lst:initial_board},
    numbers=left,
    basicstyle=\small\ttfamily,
    frame=single]
initial_board(Size, Board) :-
    length(Board, Size),        % Board tem Size linhas
    maplist(make_row(Size), Board). % cada linha tem Size colunas

make_row(Size, Row) :-
    length(Row, Size),
    maplist(=('.'), Row).       % cada celula e '.'
\end{lstlisting}

\textbf{Como funciona:} \texttt{length(Board,~Size)} cria uma lista de
\texttt{Size} vari\'{a}veis n\~{a}o instanciadas. O \texttt{maplist} aplica
\texttt{make\_row} a cada elemento, e dentro de cada linha, \texttt{maplist(=('.'),~Row)}
unifica cada c\'{e}lula com o \'{a}tomo \texttt{'.'}.

A atualiza\c{c}\~{a}o de uma c\'{e}lula \'{e} \textbf{n\~{a}o destrutiva}: o Prolog n\~{a}o
modifica estruturas existentes, mas constr\'{o}i uma nova vers\~{a}o do tabuleiro.
Esta \'{e} uma caracter\'{i}stica fundamental da programa\c{c}\~{a}o l\'{o}gica.

\begin{lstlisting}[language=Prolog,
    caption={Atualiza\c{c}\~{a}o n\~{a}o destrutiva de uma c\'{e}lula},
    label={lst:set_cell},
    numbers=left,
    basicstyle=\small\ttfamily,
    frame=single]
set_cell(Board, Row, Col, Value, NewBoard) :-
    nth0(Row, Board, OldRow, RestRows),   % separa a linha Row
    nth0(Col, OldRow, _, RestCols),       % separa o elemento Col
    nth0(Col, NewRow, Value, RestCols),   % reconstroi linha com novo valor
    nth0(Row, NewBoard, NewRow, RestRows).% reconstroi tabuleiro
\end{lstlisting}

\textbf{Analogia:} pense num fotograma de cinema. Quando um personagem
se move, n\~{a}o se apaga o fotograma anterior --- cria-se um novo fotograma
com o personagem na nova posi\c{c}\~{a}o. O Prolog funciona da mesma forma.

% ===================================================================
\section{Regras do Jogo (\texttt{rules.pl})}
\label{sec:rules}

\subsection{Valida\c{c}\~{a}o de Movimentos}

Um movimento \'{e} v\'{a}lido se a casa de destino:
\begin{enumerate}
  \item estiver a exatamente um passo (em qualquer das 8 dire\c{c}\~{o}es);
  \item estiver dentro dos limites do tabuleiro;
  \item estiver livre (\texttt{'.'}).
\end{enumerate}

\begin{lstlisting}[language=Prolog,
    caption={Valida\c{c}\~{a}o de um movimento},
    label={lst:valid_move},
    numbers=left,
    basicstyle=\small\ttfamily,
    frame=single]
valid_move(Board, FromR, FromC, ToR, ToC) :-
    board_size(Board, Size),
    Max is Size - 1,
    member(DR, [-1, 0, 1]),     % delta de linha: -1, 0 ou +1
    member(DC, [-1, 0, 1]),     % delta de coluna: -1, 0 ou +1
    \+ (DR =:= 0, DC =:= 0),   % nao pode ficar no lugar
    ToR is FromR + DR,
    ToC is FromC + DC,
    ToR >= 0, ToR =< Max,       % dentro dos limites
    ToC >= 0, ToC =< Max,
    get_cell(Board, ToR, ToC, '.'). % casa livre
\end{lstlisting}

Este predicado funciona tanto como \textbf{verificador} (quando
\texttt{ToR/ToC} j\'{a} est\~{a}o instanciados) como \textbf{gerador} (quando
est\~{a}o livres, o Prolog gera por \textit{backtracking} todas as combina\c{c}\~{o}es
de \texttt{DR/DC}). Esta dualidade \'{e} uma das caracter\'{i}sticas mais elegantes
do Prolog.

\begin{figure}[htbp]
  \centering
  \begin{minipage}[t]{0.45\textwidth}
    \centering
    \textbf{Pe\~{a}o no centro: 8 movimentos}
    \begin{verbatim}
  X X X
  X P X
  X X X
    \end{verbatim}
  \end{minipage}
  \hfill
  \begin{minipage}[t]{0.45\textwidth}
    \centering
    \textbf{Pe\~{a}o no canto: 3 movimentos}
    \begin{verbatim}
  P X .
  X X .
  . . .
    \end{verbatim}
  \end{minipage}
  \caption{Movimentos v\'{a}lidos a partir de posi\c{c}\~{o}es diferentes
           (X = destino v\'{a}lido, P = pe\~{a}o).}
  \label{fig:valid_moves}
\end{figure}

\subsection{Aplica\c{c}\~{a}o de Movimentos e Remo\c{c}\~{a}o}

\begin{lstlisting}[language=Prolog,
    caption={Aplica um movimento e remove uma casa},
    label={lst:apply_move},
    numbers=left,
    basicstyle=\small\ttfamily,
    frame=single]
% Move o peao de From para To
apply_move(Board, FromR, FromC, ToR, ToC, NewBoard) :-
    get_cell(Board, FromR, FromC, Piece),  % le o simbolo
    set_cell(Board, FromR, FromC, '.', B1),% limpa origem
    set_cell(B1, ToR, ToC, Piece, NewBoard).% coloca no destino

% Marca uma casa como permanentemente removida
apply_remove(Board, Row, Col, NewBoard) :-
    set_cell(Board, Row, Col, '#', NewBoard).
\end{lstlisting}

\subsection{Condi\c{c}\~{a}o de Fim de Jogo}

O jogo termina quando o jogador cujo turno \'{e} jogar n\~{a}o tem nenhum
movimento v\'{a}lido:

\begin{lstlisting}[language=Prolog,
    caption={Detec\c{c}\~{a}o de fim de jogo},
    label={lst:game_over},
    numbers=left,
    basicstyle=\small\ttfamily,
    frame=single]
game_over(Board, Player, P1Pos, P2Pos) :-
    player_pos(Player, P1Pos, P2Pos, CR/CC),
    \+ valid_move(Board, CR, CC, _, _).
    % \+ significa negacao: true se valid_move FALHA
\end{lstlisting}

O operador \texttt{\textbackslash+} \'{e} a \textbf{nega\c{c}\~{a}o por falha} do Prolog:
``\texttt{\textbackslash+~Goal}'' \'{e} verdade se e s\'{o} se \texttt{Goal} falha.
Aqui, \texttt{game\_over} \'{e} verdade se n\~{a}o existe nenhum movimento v\'{a}lido.

% ===================================================================
\section{Intelig\^{e}ncia Artificial (\texttt{ai.pl})}
\label{sec:ai}

\subsection{O Princ\'{i}pio Minimax}
\label{subsec:minimax}

\subsubsection{Conceito e Analogia}

O \textit{minimax} \'{e} o algoritmo cl\'{a}ssico para jogos de dois jogadores com
informa\c{c}\~{a}o perfeita~\cite{russell2020artificial}.

\textbf{Analogia:} imagine que voc\^{e} est\'{a} a jogar xadrez e quer decidir o
melhor movimento. Pensa: ``Se eu jogar aqui, o meu advers\'{a}rio vai
responder com o seu melhor movimento, e depois eu vou responder com o
meu melhor, e assim por diante.'' O \textit{minimax} formaliza exatamente
este racioc\'{i}nio.

\subsubsection{Funcionamento}

O algoritmo constr\'{o}i mentalmente uma \textbf{\'{a}rvore de jogo}:

\begin{itemize}
  \item \textbf{N\'{o}s MAX} (Jogador~1, s\'{i}mbolo \texttt{x}): escolhe o
        movimento que \textbf{maximiza} a pontua\c{c}\~{a}o.
  \item \textbf{N\'{o}s MIN} (Jogador~2, s\'{i}mbolo \texttt{o}): escolhe o
        movimento que \textbf{minimiza} a pontua\c{c}\~{a}o (prejudica o Jogador~1).
  \item \textbf{N\'{o}s folha} (profundidade~0 ou jogo terminado): a
        posi\c{c}\~{a}o \'{e} avaliada por uma fun\c{c}\~{a}o heur\'{i}stica.
\end{itemize}

\begin{figure}[htbp]
  \centering
  \begin{verbatim}
               MAX [3]
              /        \
         MIN [3]       MIN [2]
        /      \       /      \
      [3]    [12]   [8]    [2]
  \end{verbatim}
  \caption{Exemplo de \'{a}rvore minimax com 2 n\'{i}veis.
           O Jogador MAX escolhe o ramo com valor m\'{a}ximo;
           o Jogador MIN escolhe o ramo com valor m\'{i}nimo.
           O valor na raiz \'{e} 3.}
  \label{fig:minimax_tree}
\end{figure}

\begin{lstlisting}[language=Prolog,
    caption={Implementa\c{c}\~{a}o do minimax puro},
    label={lst:minimax},
    numbers=left,
    basicstyle=\small\ttfamily,
    frame=single]
% Caso base 1: profundidade esgotada -> avaliacao estatica
minimax(Board, _Player, P1Pos, P2Pos, 0, Value) :-
    !, evaluate(Board, P1Pos, P2Pos, Value).

% Caso base 2: jogador sem movimentos -> derrota terminal
minimax(Board, Player, P1Pos, P2Pos, _D, Value) :-
    player_pos(Player, P1Pos, P2Pos, CR/CC),
    \+ valid_move(Board, CR, CC, _, _), !,
    terminal_value(Player, Value).

% Caso recursivo: gera filhos e aplica max/min
minimax(Board, Player, P1Pos, P2Pos, Depth, Value) :-
    all_move_pairs(Board, Player, P1Pos, P2Pos, Pairs),
    NextD is Depth - 1,
    Other is 3 - Player,        % 1->2, 2->1
    maplist(minimax_child(Board, Player, Other,
                          P1Pos, P2Pos, NextD),
            Pairs, ChildValues),
    ( Player =:= 1
    -> max_list(ChildValues, Value)  % MAX escolhe o maior
    ;  min_list(ChildValues, Value)  % MIN escolhe o menor
    ).
\end{lstlisting}

\subsubsection{Valores Terminais}

Quando um jogador fica sem movimentos, o jogo termina:

\begin{lstlisting}[language=Prolog,
    caption={Valores terminais},
    label={lst:terminal},
    numbers=left,
    basicstyle=\small\ttfamily,
    frame=single]
terminal_value(1, -1000). % Jogador 1 preso -> derrota (negativo)
terminal_value(2,  1000). % Jogador 2 preso -> vitoria de P1 (positivo)
\end{lstlisting}

Os valores $\pm1000$ s\~{a}o intencionalmente muito maiores do que qualquer
valor de avalia\c{c}\~{a}o heur\'{i}stica, garantindo que uma vit\'{o}ria ou derrota
real \'{e} sempre preferida a um estado com boa heur\'{i}stica.

% -------------------------------------------------------------------
\subsection{Algoritmo Alfa-Beta}
\label{subsec:alphabeta}

\subsubsection{Conceito e Analogia}

O \textit{minimax} \'{e} correto mas lento: explora todos os n\'{o}s da \'{a}rvore.
O \textit{alfa-beta} \'{e} uma otimiza\c{c}\~{a}o que descarta ramos que \textbf{nunca
podem influenciar o resultado}, sem alterar o valor final calculado.

\textbf{Analogia:} imagine que est\'{a} a selecionar o melhor restaurante numa
lista. J\'{a} visitou um restaurante com classifica\c{c}\~{a}o 8/10. Ao visitar o
segundo restaurante, descobre que tem um prato horrível de 3/10. Sem sequer
provar o resto da ementa, sabe que este restaurante nunca ser\'{a} a sua
primeira escolha. \textbf{Pode ignorar o resto da visita} e passar ao
pr\'{o}ximo. O alfa-beta faz exatamente o mesmo com ramos da \'{a}rvore de jogo.

\subsubsection{Os Par\^{a}metros $\alpha$ e $\beta$}

\begin{itemize}
  \item $\alpha$ (\textit{alpha}) --- o melhor valor que o Jogador MAX j\'{a}
        garantiu. Come\c{c}a em $-\infty$.
  \item $\beta$ (\textit{beta}) --- o melhor valor que o Jogador MIN j\'{a}
        garantiu. Come\c{c}a em $+\infty$.
  \item \textbf{Corte beta} (em n\'{o} MAX): se o valor atual $\geq \beta$,
        o Jogador MIN nunca escolheria este ramo. Pode parar.
  \item \textbf{Corte alfa} (em n\'{o} MIN): se o valor atual $\leq \alpha$,
        o Jogador MAX nunca escolheria este ramo. Pode parar.
\end{itemize}

\begin{figure}[htbp]
  \centering
  \begin{verbatim}
                MAX [3]
               /        \
          MIN [3]       MIN [?]
         /      \       /      \
       [3]    [12]   [2]   CORTADO
                     ^
          Alpha = 3, este MIN ja tem 2.
          Como 2 < 3 (alpha), MAX nunca
          escolheria este caminho.
          Os filhos restantes sao cortados.
  \end{verbatim}
  \caption{Corte alfa: o Jogador MAX j\'{a} tem garantido o valor~3
           (ramo esquerdo). No ramo direito, o MIN encontrou 2,
           que \'{e} pior para MAX. Os restantes filhos do ramo
           direito s\~{a}o descartados.}
  \label{fig:alphabeta_cutoff}
\end{figure}

\subsubsection{Implementa\c{c}\~{a}o}

\begin{lstlisting}[language=Prolog,
    caption={Pesquisa alfa-beta --- caso recursivo},
    label={lst:alphabeta},
    numbers=left,
    basicstyle=\small\ttfamily,
    frame=single]
alphabeta(Board, Player, P1Pos, P2Pos, Depth, Alpha, Beta,
          Value, BM, BR) :-
    all_move_pairs(Board, Player, P1Pos, P2Pos, Pairs),
    order_pairs(Pairs, Board, Player, P1Pos, P2Pos, Ordered),
    NextDepth is Depth - 1,
    Other is 3 - Player,
    ( Player =:= 1
    ->  ab_max(Ordered, Board, Other, P1Pos, P2Pos,
               Player, NextDepth, Alpha, Beta,
               -10000, none, none, Value, BM, BR)
    ;   ab_min(Ordered, Board, Other, P1Pos, P2Pos,
               Player, NextDepth, Alpha, Beta,
               10000, none, none, Value, BM, BR)
    ).

% No MAX: actualiza alpha e corta se alpha >= beta
ab_max([move(MPos,RPos)|Rest], Board, Other, P1Pos, P2Pos,
       CurPlayer, Depth, Alpha, Beta, BestVal, BestM, BestR,
       Val, BM, BR) :-
    apply_pair(Board, CurPlayer, P1Pos, P2Pos, MPos, RPos,
               NewBoard, NewP1, NewP2),
    alphabeta(NewBoard, Other, NewP1, NewP2, Depth,
              Alpha, Beta, ChildVal, _, _),
    ( (BestM = none ; ChildVal > BestVal)
    ->  NewBest = ChildVal, NewBM = MPos, NewBR = RPos
    ;   NewBest = BestVal, NewBM = BestM, NewBR = BestR
    ),
    NewAlpha is max(Alpha, NewBest),
    ( NewAlpha >= Beta
    ->  Val = NewBest, BM = NewBM, BR = NewBR  % corte beta
    ;   ab_max(Rest, Board, Other, P1Pos, P2Pos,
               CurPlayer, Depth, NewAlpha, Beta,
               NewBest, NewBM, NewBR, Val, BM, BR)
    ).
\end{lstlisting}

\subsubsection{Ganho de Desempenho}

O corte alfa-beta reduz drasticamente o n\'{u}mero de n\'{o}s avaliados~\cite{knuth1975analysis}:

\begin{table}[htbp]
  \caption{Complexidade: minimax vs.\ alfa-beta}
  \label{tab:complexity}
  \centering
  \begin{tabular}{lll}
    \toprule
    \textbf{Algoritmo} & \textbf{Complexidade} & \textbf{N\'{o}s (b=8, d=4)} \\
    \midrule
    Minimax puro    & $O(b^d)$      & $8^4 = 4{.}096$ \\
    Alfa-beta ideal & $O(b^{d/2})$  & $8^2 = 64$ \\
    \bottomrule
  \end{tabular}
\end{table}

Na pr\'{a}tica, considerando que cada movimento tem em m\'{e}dia~8 destinos
poss\'{i}veis e 34 casas remov\'{i}veis num tabuleiro 6{\texttimes}6 inicial,
o fator de ramifica\c{c}\~{a}o real \'{e} $\approx 102$ (ver
Cap\'{i}tulo~\ref{chap:eval}). O alfa-beta com boa ordena\c{c}\~{a}o reduz este
fator para $\approx \sqrt{102} \approx 10$ por n\'{i}vel.

% -------------------------------------------------------------------
\subsection{Ordena\c{c}\~{a}o de Movimentos}
\label{subsec:ordering}

A efic\'{a}cia do corte alfa-beta depende fortemente da ordem em que os
movimentos s\~{a}o avaliados. Avaliar primeiro os movimentos mais
prometedores maximiza os cortes.

\textbf{Heur\'{i}stica implementada:}
\begin{itemize}
  \item \textbf{Movimentos centrais primeiro}: casas perto do centro do
        tabuleiro oferecem mais op\c{c}\~{o}es futuras (maior mobilidade).
  \item \textbf{Remo\c{c}\~{o}es pr\'{o}ximas do advers\'{a}rio primeiro}: remover casas
        adjacentes ao advers\'{a}rio reduz mais rapidamente a sua mobilidade.
\end{itemize}

\begin{lstlisting}[language=Prolog,
    caption={Fun\c{c}\~{a}o de pontua\c{c}\~{a}o para ordena\c{c}\~{a}o de movimentos},
    label={lst:ordering},
    numbers=left,
    basicstyle=\small\ttfamily,
    frame=single]
score_pair(Center, OppR, OppC,
           move(MR/MC, RR/RC),
           Score-move(MR/MC, RR/RC)) :-
    MDR is abs(MR - Center), MDC is abs(MC - Center),
    MoveScore is MDR + MDC,    % menor = mais central
    RDR is abs(RR - OppR), RDC is abs(RC - OppC),
    RemScore is RDR + RDC,     % menor = mais perto do adversario
    Score is MoveScore * 3 - RemScore. % peso: movimento > remocao
\end{lstlisting}

% -------------------------------------------------------------------
\subsection{Fun\c{c}\~{a}o de Avalia\c{c}\~{a}o}
\label{subsec:eval_func}

Quando a pesquisa atinge a profundidade m\'{a}xima (sem o jogo ter
terminado), \'{e} necess\'{a}rio avaliar heuristicamente a posi\c{c}\~{a}o atual.

\textbf{Dois crit\'{e}rios foram escolhidos:}

\begin{enumerate}
  \item \textbf{Mobilidade} --- o n\'{u}mero de movimentos legais dispon\'{i}veis.
        Um jogador com mais movimentos tem mais liberdade e est\'{a} mais longe
        de ser isolado. Este crit\'{e}rio tem peso~10.
  \item \textbf{Centralidade} --- qu\~{a}o pr\'{o}ximo est\'{a} o pe\~{a}o do centro
        do tabuleiro. O centro oferece mais dire\c{c}\~{o}es de fuga.
        Este crit\'{e}rio serve para desempatar.
\end{enumerate}

\begin{lstlisting}[language=Prolog,
    caption={Fun\c{c}\~{a}o de avalia\c{c}\~{a}o heur\'{i}stica},
    label={lst:evaluate},
    numbers=left,
    basicstyle=\small\ttfamily,
    frame=single]
evaluate(Board, P1Pos, P2Pos, Value) :-
    P1Pos = P1R/P1C, P2Pos = P2R/P2C,
    valid_moves(Board, P1R, P1C, M1), length(M1, S1), % mobilidade P1
    valid_moves(Board, P2R, P2C, M2), length(M2, S2), % mobilidade P2
    board_size(Board, Size),
    Center is (Size - 1) / 2.0,
    centrality(P1R, P1C, Center, Cent1),
    centrality(P2R, P2C, Center, Cent2),
    Value is (S1 - S2) * 10 + (Cent1 - Cent2).
    % positivo = bom para P1; negativo = bom para P2
\end{lstlisting}

\textbf{Exemplo pr\'{a}tico:} num tabuleiro 3{\texttimes}3 com P1 no canto (0,0)
e P2 no centro (1,1), o c\'{a}lculo \'{e}:
\begin{itemize}
  \item Mobilidade P1 = 2 movimentos (confinado ao canto)
  \item Mobilidade P2 = 7 movimentos (centro, m\'{a}xima liberdade)
  \item Diferen\c{c}a de mobilidade: $(2-7) \times 10 = -50$
  \item Centralidade: $(0+0) - (1+1) = -2$ (P2 est\'{a} mais ao centro)
  \item Pontua\c{c}\~{a}o total: $-50 + (-2) = -52$
\end{itemize}
O valor negativo indica que P2 est\'{a} em clara vantagem --- o que \'{e} correto:
P2 no centro tem muito mais liberdade.

\textbf{Justifica\c{c}\~{a}o da escolha:} a mobilidade \'{e} o indicador mais direto
de quanto um jogador est\'{a} pr\'{o}ximo de ser isolado. A centralidade captura
a tend\^{e}ncia estrat\'{e}gica de manter o pe\~{a}o no centro do tabuleiro.

% ===================================================================
\section{Interface de Jogo (\texttt{game.pl})}
\label{sec:game_loop}

\subsection{Fluxo do Programa}

O jogo segue o seguinte fluxo:

\begin{enumerate}
  \item \texttt{play/0} --- exibe o menu principal.
  \item \texttt{main\_menu/0} --- permite escolher o modo de jogo
        (humano-humano, humano-IA, IA-IA).
  \item \texttt{setup\_game/2} --- configura o tamanho do tabuleiro e a
        profundidade de pesquisa da IA.
  \item \texttt{start\_game/4} --- cria o tabuleiro inicial com os pe\~{o}es
        nos cantos opostos.
  \item \texttt{game\_loop/6} --- ciclo principal: exibe o tabuleiro,
        verifica fim de jogo, executa o turno do jogador corrente.
\end{enumerate}

\begin{lstlisting}[language=Prolog,
    caption={Ciclo principal do jogo},
    label={lst:game_loop},
    numbers=left,
    basicstyle=\small\ttfamily,
    frame=single]
game_loop(Board, Turn, P1Pos, P2Pos, Mode, Depth) :-
    display_game(Board, Turn),
    ( game_over(Board, Turn, P1Pos, P2Pos)
    ->  Winner is 3 - Turn,  % o adversario venceu
        announce_winner(Winner, Board),
        play_again
    ;   execute_turn(Board, Turn, P1Pos, P2Pos,
                     Mode, Depth,
                     NewBoard, NewP1, NewP2),
        NextTurn is 3 - Turn,  % alterna: 1->2, 2->1
        game_loop(NewBoard, NextTurn, NewP1, NewP2,
                  Mode, Depth)
    ).
\end{lstlisting}

\subsection{Turno Humano vs.\ Turno IA}

Quando o turno pertence a um jogador humano, o programa pede a
entrada no formato \texttt{Linha/Coluna.} (com ponto final, pois o
SWI-Prolog interpreta a entrada como um termo Prolog completo):

\begin{verbatim}
Move (linha/coluna, ex: 1/2): 1/1.
Remove uma casa (linha/coluna): 2/2.
\end{verbatim}

Quando o turno pertence \`{a} IA, o predicado \texttt{best\_move/7}
\'{e} invocado automaticamente, que executa o alfa-beta e devolve o
par (movimento, remo\c{c}\~{a}o) \'{o}timo.

% ===================================================================
\section{Utiliza\c{c}\~{a}o de Intelig\^{e}ncia Artificial no Desenvolvimento}
\label{sec:ai_usage}

Ao longo do desenvolvimento deste projeto, a IA generativa
(nomeadamente o modelo Claude da Anthropic) foi utilizada de forma
extensa e integrada em todas as fases do trabalho. A sua utiliza\c{c}\~{a}o
n\~{a}o se limitou a gerar c\'{o}digo, mas abrangeu todo o processo de
aprendizagem e desenvolvimento:

\begin{enumerate}
  \item \textbf{Aprendizagem de conceitos:} a IA foi consultada para
        compreender os fundamentos do Prolog (unifica\c{c}\~{a}o,
        \textit{backtracking}, listas, predicados embutidos) atrav\'{e}s
        de explica\c{c}\~{o}es, exemplos e analogias.

  \item \textbf{Compreens\~{a}o dos algoritmos:} o princ\'{i}pio
        \textit{minimax} e o algoritmo \textit{alfa-beta} foram
        explorados com a assist\^{e}ncia da IA, que forneceu
        explica\c{c}\~{o}es passo a passo, diagramas de \'{a}rvores e
        an\'{a}lise de complexidade.

  \item \textbf{Explora\c{c}\~{a}o de alternativas:} foram discutidas
        diferentes abordagens para a representa\c{c}\~{a}o do tabuleiro
        (lista de listas vs.\ tabela de hash), para a fun\c{c}\~{a}o de
        avalia\c{c}\~{a}o (mobilidade, centralidade, dist\^{a}ncia entre
        pe\~{o}es) e para a ordena\c{c}\~{a}o de movimentos.

  \item \textbf{Escrita e depura\c{c}\~{a}o de c\'{o}digo:} a IA auxiliou na
        escrita de predicados, na corre\c{c}\~{a}o de erros l\'{o}gicos
        (incluindo um \textit{bug} na acumula\c{c}\~{a}o de valores no
        alfa-beta), e na formula\c{c}\~{a}o dos 33 testes unit\'{a}rios.

  \item \textbf{Elabora\c{c}\~{a}o do relat\'{o}rio:} a IA foi utilizada para
        auxiliar na reda\c{c}\~{a}o, estrutura\c{c}\~{a}o e revis\~{a}o deste
        documento, com o objetivo de tornar as explica\c{c}\~{o}es claras
        e acess\'{i}veis.
\end{enumerate}

Em todos os casos, o uso da IA foi orientado pelos objetivos
did\'{a}ticos da unidade curricular: a IA serviu como ferramenta de
aprendizagem e produtividade, e n\~{a}o como substituto da
compreens\~{a}o dos conceitos.
